

\section{Basic data and properties}
\label{sec:6.4}

The \textsl{rank  of the kernel}, $r$, satisfies $\pi N_0=\frac{1}{r}\,u$.\\
A related quantity is $r'=1-\frac{1}{r}$. Then $\pi N=r'u$, $r'\dg  u.=N\dg  \pi.$, and $\Omega N=r'J$.

For an idempotent $e$, the $i^{th}$ component of $uee^*$ is the size of the block containing $i$.
For an element of $k$ having the same partition as $e$,
the squared norm, $\|uk\|^2=ukk^*u^\dagger=uee^*\dg u.$, is the sum of the squares of the block sizes. 
Thus, the \textsl{average block vector\/}, $b$,  
defined by $b=uN_0$, has components $b_i=\avg B_i.$ where $B_i$ is the random variable
giving the size of the block containing $i$ running through the partitions of $\cK$. 
Note that $ukk^*u^\dagger=\tr k^*Jk$, and averaging yields
$\displaystyle\tau=\avg\sum_{i=1}^rB_i^2.$, the expected sum of squares of the block sizes running over all partitions.

We have as well $\tau=\tr JN_0$. A quantity related to $\tau$ is $\tau'=n^2-\tau$. Then we have $\tau'=\tr JN=\tr JM$.

Note that $uM=\tau'\pi$, $M\dg  u.=\tau'\dg  \pi.$ and $JM=\tau' \Omega$.

In the next section we will prove the result of BM-F that says $\pi w e=(1/r)\rho(e)=\pi e$ for any word $w$ 
in $\cS$. Consequently, $\pi wN=\pi N=r'u$ for any $w\in\cS$.
